\subsection{Start Spark Shell}

\par
To allow the Spark containers to access the Postgres JDBC driver, I mounted the
\texttt{jars/} folder into each Spark service via a Docker volume. This makes the
jar available at \texttt{/opt/spark/jars/extra/} inside the containers:

\begin{figure}[H]
    \centering
    \includegraphics[width=0.9\textwidth]{2_6_postgres_jar.png}
    \caption{Volume mount added to Spark services in \texttt{docker-compose.yml} exposing the Postgres JDBC jar}
\end{figure}

\par
After updating the compose file, I restarted the Spark services to apply the new
volume configuration:

\begin{figure}[H]
    \centering
    \includegraphics[width=0.9\textwidth]{2_6_containers_restarted.png}
    \caption{Spark master and worker containers restarted with the new volume mount}
\end{figure}

\par
I then opened a shell inside the \texttt{spark-master} container and launched
\texttt{spark-shell} pointing to the Standalone master, with 2 executors matching
the 2 workers, the Postgres JDBC jar, and the MongoDB Spark connector package.
The Ivy cache was redirected to \texttt{/tmp/.ivy2} to avoid permission issues with
the default home directory:

\begin{lstlisting}[language=bash]
bin/spark-shell \
  --master spark://spark-master:7077 \
  --num-executors 2 \
  --jars /opt/spark/jars/extra/postgresql-42.7.3.jar \
  --packages org.mongodb.spark:mongo-spark-connector_2.12:10.3.0 \
  --conf spark.driver.extraJavaOptions="-Divy.home=/tmp/.ivy2 -Divy.cache.dir=/tmp/.ivy2/cache"
\end{lstlisting}

\begin{figure}[H]
    \centering
    \includegraphics[width=0.9\textwidth]{2_6_spark_shell_started.png}
    \caption{Spark shell successfully started with the \texttt{scala>} prompt}
\end{figure}

\par
To verify the running application, I attempted to access the Spark Application WebUI
at \texttt{http://localhost:4040}. The port was not initially exposed, so I added it
to the \texttt{spark-master} service in \texttt{docker-compose.yml}:

\begin{figure}[H]
    \centering
    \includegraphics[width=0.9\textwidth]{2_6_expose_4040.png}
    \caption{Port 4040 added to \texttt{spark-master} in \texttt{docker-compose.yml}}
\end{figure}

\par
A \texttt{docker compose restart} on the master alone was not sufficient to bind
the new port --- a full \texttt{docker compose down} and \texttt{up} was required:

\begin{figure}[H]
    \centering
    \includegraphics[width=0.9\textwidth]{2_6_restart_master.png}
    \caption{Restarting only the master --- port 4040 still not bound}
\end{figure}

\begin{figure}[H]
    \centering
    \includegraphics[width=0.9\textwidth]{2_6_full_restart.png}
    \caption{Full \texttt{docker compose down} and \texttt{up} successfully binds port 4040}
\end{figure}

\par
With the shell running and port 4040 exposed, the Spark Application WebUI became
accessible at \texttt{http://localhost:4040}, confirming the application was running
correctly:

\begin{figure}[H]
    \centering
    \includegraphics[width=0.9\textwidth]{2_6_spark_app_webui.png}
    \caption{Spark Application WebUI at \texttt{http://localhost:4040} showing the running shell application}
\end{figure}
